% ==============================================================================
% РАЗДЕЛ 2.3 (DISS-005-B1393 .. DISS-005-B1444)
% Совершенствование алгоритмов функционирования восстановления нормального режима
% ==============================================================================

% DISS-005-B1393
\section[Совершенствование алгоритмов функционирования восстановления нормального режима]{Совершенствование алгоритмов функционирования восстановления нормального режима\commentnote{DISS-005-C0096}{Алексей Евдаков [2]}{Требуется скорректировать текст, чтобы не было копипастом со статьи}\commentnote{DISS-005-C0097}{Евдаков Алексей}{Взято из статьи на Бенардосовские чтения 2025\\ \\ СОВЕРШЕНСТВОВАНИЕ АЛГОРИТМОВ ФУНКЦИОНИРОВАНИЯ ВОССТАНОВЛЕНИЯ НОРМАЛЬНОГО РЕЖИМА}}
\label{sec:ch02_vnr_algorithms}

% DISS-005-B1394
С целью повышения надежности электроснабжения, сокращения времени перерыва питания потребителей и упрощения работы оперативного персонала [40, 41] на современных подстанциях предусматривается наличие нескольких независимых источников питания. Для автоматического переключения между ними применяются функции автоматического ввода резерва (АВР) и восстановления нормального режима (ВНР).\par\vspace{0.2cm}

% DISS-005-B1395
Алгоритм АВР обеспечивает переключение питания секции шин, оказавшейся без питания (например, вследствие внешнего короткого замыкания с последующим отключением питающей линии электропередачи), на резервный источник питания, работа которого не была нарушена аварийным процессом (Рисунок~2.15.а). Восстановление исходной (нормальной) схемы электроснабжения осуществляется либо вручную оперативным персоналом, либо автоматически посредством функции ВНР (Рисунок~2.15.б).\par\vspace{0.2cm}

% DISS-005-B1396
Вследствие происходящего цифрового прогресса и возможности реализации сложных логических схем на микропроцессорных устройствах, были созданы быстродействующие АВР (БАВР). Данные устройства способные обеспечивать переключения на резервный источник питания за десятки миллисекунд, что критически важно для технологических процессов в ряде отраслей промышленности. Современные терминалы БАВР, как правило, реализуют и функцию ВНР.\par\vspace{0.2cm}

% DISS-005-B1397
За последние десятилетия на энергообъектах внедрено значительное количество терминалов БАВР различных производителей, реализующих и функцию ВНР. Алгоритмы функционирования данных устройств, включая процессы ликвидации аварийных режимов и последующего восстановления нормальной схемы, претерпели существенные изменения. Ряд таких усовершенствований функционирования ВНР рассматриваются в данной статье далее.\par\vspace{0.2cm}

\begin{figure}[htbp]
  \centering
  % DISS-005-B1398 (DISS-005-FIG-078)
  \includegraphics[width=0.75\textwidth]{image81.png}\par\vspace{1mm}
  % DISS-005-B1399
  {\small а)}\par\vspace{3mm}
  % DISS-005-B1400 (DISS-005-FIG-079)
  \includegraphics[width=0.75\textwidth]{image82.png}\par\vspace{1mm}
  % DISS-005-B1401
  {\small б)}\par\vspace{2mm}
  % DISS-005-B1402
  \caption[Логика подачи команд автоматикой при возникновении аварии (АВР) и его последующего восстановления (ВНР)]{Логика подачи команд автоматикой при возникновении аварии (АВР) и его последующего восстановления (ВНР) на примере первой секции: а) логика АВР б) логика работы ВНР}
  \label{fig:ch02_logic_avr_vnr}
\end{figure}

% DISS-005-B1403

% DISS-005-B1404
\subsection{Существующие требования к ВНР}
\label{sec:ch02_vnr_existing_requirements}

% DISS-005-B1405
Согласно требованиям нормативных документам [42] устройства ВНР должны обеспечивать контроль восстановления напряжения, отсутствие пуска или срабатывания токовых защит на вводах, а также отстройку от трёхкратного автоматического повторного включения (АПВ) на линиях.\par\vspace{0.2cm}

% DISS-005-B1406
Однако перечисленных условий может быть недостаточно для предотвращения некорректных переключений ВНР в ряде сложных режимов, которые могут привести к аварийным ситуациям или недопустимым перегрузкам оборудования, вызывающим последующие отключения.\par\vspace{0.2cm}

% DISS-005-B1407
\subsection{Последовательный и параллельный ВНР}
\label{sec:ch02_vnr_series_parallel}

% DISS-005-B1408
На многих промышленных предприятиях применяется алгоритм параллельного ВНР, предусматривающий кратковременную параллельную работу источников питания. Такой подход позволяет обеспечить непрерывность питания для нагрузки в процессе переключений, однако, в отличие от алгоритма последовательного ВНР, он несет в себе риски развития аварийных ситуаций при определенных условиях, рассмотренных далее.\par\vspace{0.2cm}

% DISS-005-B1409
\subsection{Блокировка по однофазным замыканиям на землю}
\label{sec:ch02_vnr_blocking_ground_fault}

% DISS-005-B1410
Одним из первых усовершенствований, внедренных многими производителями [43, 44, 45, 46], стала блокировка ВНР при возникновении однофазных замыканий на землю (ОЗЗ). Применяются различные варианты реализации данной блокировки\par\vspace{0.2cm}

% DISS-005-B1411
Блокировка при наличии ОЗЗ на обеих секциях шин. Данный алгоритм предотвращает возникновение междуфазного короткого замыкания (КЗ) при включении на параллельную работу секций с ОЗЗ на разных фазах.\par\vspace{0.2cm}

% DISS-005-B1412
Блокировка при наличии ОЗЗ на любой из секций шин. Этот вариант дополнительно снижает риск возникновения ОЗЗ на ранее неповрежденной секции или перерастания существующего ОЗЗ в более тяжелый вид повреждения (например, двойное замыкание на землю).\par\vspace{0.2cm}

% DISS-005-B1413
Данная блокировка является стандартной для современных терминалов БАВР с функцией ВНР и позволяет предотвратить развитие очевидных аварийных режимов при ОЗЗ.\par\vspace{0.2cm}

% DISS-005-B1414
\subsection{Блокировка при превышении угла}
\label{sec:ch02_vnr_blocking_angle_exceedance}

% DISS-005-B1415
С целью ускорения срабатывания АВР (особенно при наличии мощной синхронной нагрузки, замедляющей снижение напряжения) в ряде терминалов БАВР, помимо контроля снижения напряжения, был внедрен пусковой орган по углу между векторами напряжений источников. Стандартные уставки срабатывания по углу для БАВР составляют 10-15$^\circ$. Превышение этого значения обычно свидетельствует о наступлении аварийного режима.\par\vspace{0.2cm}

% DISS-005-B1416
Однако при проведении ремонтных работ, на ряде объектов угол между секциями может составлять 15-25$^\circ$. Такое значение угла выходит за допустимые диапазоны, и даже применение адаптивных уставок доступных у некоторых производителей [43] не всегда исключает ложное срабатывание БАВР в подобных режимах (Рисунок~2.16).\par\vspace{0.2cm}

% DISS-005-B1417
До внедрения дополнительных блокировок в таких ситуациях наблюдались циклические срабатывания АВР--ВНР. Это происходило потому, что напряжение на вводах по модулю оставалось в пределах нормы, и после отработки выдержки времени готовности ВНР происходило восстановление нормальной схемы, что приводило к повторному срабатыванию БАВР по углу. Эксплуатирующему персоналу приходилось временно выводить из работы функции АВР и ВНР на период таких режимов.\par\vspace{0.2cm}

% DISS-005-B1418
Для решения данной проблемы были внедрены следующие две блокировки ВНР:\par\vspace{0.2cm}

% DISS-005-B1419
Блокировка ВНР по углу. Предотвращает срабатывание ВНР, если угол превышает заданное допустимое значение (значение должно быть меньше уставки пуска по углу).\par\vspace{0.2cm}

% DISS-005-B1420
Блокировка ВНР по числу срабатываний за определенный интервал времени. Предотвращает циклические переключения после их наступления.\par\vspace{0.2cm}

\begin{figure}[htbp]
  \centering
  \begin{minipage}{\textwidth}
    \centering
    % DISS-005-B1421 (DISS-005-FIG-080)
    \includegraphics[width=0.95\textwidth]{image83.png}\commentnote{DISS-005-C0098}{Алексей Евдаков [2]}{Найти осциллограмму позже}
    % DISS-005-B1422
    \caption{Срабатывание БАВР по превышению угла при проведении ремонтных работ.}
    \label{fig:ch02_bavr_angle_exceedance_repair}
  \end{minipage}
\end{figure}

% DISS-005-B1423

% DISS-005-B1424
\subsection{Блокировка по частоте}
\label{sec:ch02_vnr_blocking_frequency}

% DISS-005-B1425
Согласно ПУЭ [47] потребители первой категории должны обеспечиваться минимум двумя независимыми источниками питания. Однако в ряде случаев, в качестве такого независимого источника выступает собственная станция, и следовательно~--- частота и угол между секциями могут не совпадать. И если синхронизма по углу можно добиться, то существенная разность по частотам приведёт к тому, что данного условия будет недостаточно, так как двигательная нагрузка обладает большой инерцией и не может мгновенно разогнаться до 50~Гц. Это приводит к большим уравнительным токам и недопустимым электромеханическим переходным процессам.\par\vspace{0.2cm}

% DISS-005-B1426
Рассмотрим показательный случай несинхронного срабатывания ВНР, произошедший на одном из крупных промышленных предприятий с собственной генерацией. В результате внешнего короткого замыкания произошло отключение одной из питающих высоковольтных линий (ВЛ 220~кВ). Устройства БАВР на подстанциях предприятия отработали штатно, переведя питание части потребителей на резервные источники, включая шины собственной генерации.\par\vspace{0.2cm}

% DISS-005-B1427
Данное событие привело к изменению режима работы локальной энергосистемы: часть оборудования предприятия начала работать изолированно от внешней сети, что сопровождалось снижением частоты (в рассматриваемом случае до 49,2~Гц) и потерей синхронизма с энергосистемой.\par\vspace{0.2cm}

% DISS-005-B1428
После восстановления напряжения на отключившейся ВЛ 220~кВ инициировалась работа автоматики ВНР на главной понизительной подстанции (ГПП) предприятия. В момент попытки параллельного включения основного и резервного источников (внешняя сеть и собственная генерация) угол между их напряжениями достигал значительной величины (порядка 140$^\circ$), а частоты существенно различались. Это вызвало большие уравнительные токи (пиковое значение достигало 16,2~кА, действующее значение~--- до 9,7~кА). Приведён пример одной из подстанций на которой бросок тока составил 7,5~кА (Рисунок~2.17).\par\vspace{0.2cm}

\begin{figure}[htbp]
  \centering
  \begin{minipage}{\textwidth}
    \centering
    % DISS-005-B1429 (DISS-005-FIG-081)
    \includegraphics[width=0.95\textwidth]{image84.png}\commentnote{DISS-005-C0099}{Алексей Евдаков [2]}{Найти осциллограмму позже}
    % DISS-005-B1430
    \caption{Срабатывание ВНР без контроля угла и очерёдности функционирования.}
    \label{fig:ch02_vnr_no_angle_order_control}
  \end{minipage}
\end{figure}

% DISS-005-B1431
Возникшие сложные переходные процессы в электрической сети ГПП повлекли за собой излишнее срабатывание релейной защиты с последующим отключением второго силового трансформатора ГПП. Это привело к полному исчезновению напряжения на части распределительных устройств и аварийному останову подключенного к ним оборудования.\par\vspace{0.2cm}

% DISS-005-B1432
Предотвратить развитие описанной аварийной ситуации, приведшей к значительным технологическим нарушениям, могли бы следующие меры, реализуемые в ряде современных устройств:\par\vspace{0.2cm}

% DISS-005-B1433
Блокировка по углу. Благодаря тому, что зона допустимого срабатывания равная 30$^\circ$ (например, при уставке пуска по углу 15$^\circ$) проходится в данном событии приблизительно за 0,11~с, то выдержка на готовность ВНР равная 10~с не успела бы набраться и не позволила бы сработать функции ВНР\par\vspace{0.2cm}

% DISS-005-B1434
Блокировка по разности частот. При отклонении частоты между секциями на значение больше допустимого, блокируется работа автоматики ВНР на длительное время.\par\vspace{0.2cm}

% DISS-005-B1435
Использование функции контроля синхронизма (ANSI 25). Применение полноценного алгоритма проверки условий синхронизма (нахождение модулей, частот и взаимного угла напряжений в допустимом диапазоне) перед выдачей команды на выключатели для параллельной работы в цикле ВНР.\par\vspace{0.2cm}

% DISS-005-B1436
\subsection{Уменьшение одновременного срабатывания ВНР}
\label{sec:ch02_vnr_reduce_simultaneous}

% DISS-005-B1437
При наличии на одном объекте (например, крупном промышленном предприятии) множества терминалов БАВР возрастает вероятность их одновременного срабатывания, что может приводить к значительным суммарным броскам тока при восстановлении нормальной схемы.\par\vspace{0.2cm}

% DISS-005-B1438
Если при функционировании логики БАВР почти одновременная работа на разных присоединениях может быть обусловлена требованиями минимизации перерыва питания ответственных потребителей, то для функции ВНР синхронное срабатывание множества устройств не является необходимым и может быть нежелательным.\par\vspace{0.2cm}

% DISS-005-B1439
В ситуации, подобной рассмотренной выше (отключение внешней питающей линии главной понизительной подстанции), срабатывание БАВР затрагивает значительное число присоединений на нижестоящих уровнях напряжения. После восстановления напряжения на основной питающей линии на всех задействованных устройствах запускается отсчет выдержки времени готовности ВНР (Рисунок~2.17). Зачастую на объекте устанавливается единая стандартная выдержка времени (например, 10~с) и одинаковое время параллельной работы (например, 0,2~с). Это приводит к практически одновременному срабатыванию множества ВНР, суммированию токов параллельной работы и, как следствие, к значительным перетокам мощности между источниками. Такое переключение, особенно при наличии других неблагоприятных факторов (например, начального расхождения угла между источниками), может вызвать превышение допустимых токов коммутационной аппаратуры или глубокие просадки напряжения в сети предприятия.\par\vspace{0.2cm}

% DISS-005-B1440
Для снижения негативного влияния одновременной работы ВНР рекомендуется в пределах одного энергообъекта применять селективность по времени срабатывания ВНР путем задания различных выдержек времени. Например, вместо единой уставки 10~с можно использовать диапазон уставок от 9 до 11~с (или шире, в зависимости от объекта) с шагом, равным или несколько превышающим время параллельной работы (например, 0,2-0,5~с), распределяя их между различными группами устройств ВНР на объекте. Это обеспечит разновременность переключений и снизит пиковые нагрузки на сеть и оборудование.:\par\vspace{0.2cm}

% DISS-005-B1441
\subsection{Выводы о совершенствование алгоритмов функционирования восстановления нормального режима}
\label{sec:ch02_vnr_conclusions}

% DISS-005-B1442
Представлены современные модернизации логики функционирования параллельного восстановления нормального режима для исключения включения на короткие замыкания и возникновения больших перетоков мощности, приводящих к аварийным отключениям, а также для уменьшения возникающих перегрузок во время переключений.\par\vspace{0.2cm}

% DISS-005-B1443
Рассмотрены актуальные усовершенствования алгоритмов функции параллельного ВНР, направленные на предотвращение некорректных срабатываний (в том числе включений на КЗ или несинхронную параллельную работу), способных вызвать значительные перетоки мощности и привести к аварийным отключениям. Внедрение блокировок по ОЗЗ, по углу между напряжениями источников, по разности частот, а также использование полноценного контроля синхронизма повышают надежность функционирования ВНР в сложных режимах. Кроме того, предложены способы снижения токовых нагрузок и перетоков мощности в процессе выполнения переключений ВНР за счет введения временной селективности между различными устройствами ВНР на одном объекте. Реализация описанных подходов позволяет повысить надежность электроснабжения потребителей и предотвратить развитие аварийных ситуаций, связанных с работой автоматики ВНР.\par\vspace{0.2cm}

% DISS-005-B1444

